\section{Periodic modules and a nonpolarization class}
\label{per:section}

We now use the ordinary mirror only at $t=0$.  The aim is to interpret
the four graded modules $V_{\rm s}^r$ as section modules of coherent
sheaves and then to detect a characteristic class by Euler pairings.
The new Lagrangian need not be an object of the ordinary mirror
category: its module will be extended from the directed category by
the repeated graph copies.

Retain the enlarged coefficient field $\kappa$ and the ordinary line
model $\mathcal B$ of Proposition~\ref{mirror:simultaneous}.  Its objects
have the graph labels, and its split-closed pretriangulated hull is
$\Perf(Y_\kappa)$.  The layer-zero labels $G_k$ and $C$ correspond to
ordinary, unshifted lines $L_k'$ and $P_C$, respectively.  Here $L_k'$
has numerical type $H_Y^{-h_0k}$.  The common shift has already been
removed in the construction of $\mathcal B$.

\subsection{Periodic module conventions}

For an integer-graded complex $K$, define its direct-sum fold by
\begin{equation}
 \operatorname{Fold}(K)^r
   =\bigoplus_{n\equiv r\;({\rm mod}\,4)}K^n,
 \qquad r\in\Z/4.
 \label{per:fold}
\end{equation}
Fold $\mathcal B$ in this way and denote the result by
$\overline{\mathcal B}$.  Fold $\mathcal D_\kappa$ and adjoin formal
strict units, obtaining $\overline{\mathcal D}$.  The nonunital ordered
map of Proposition~\ref{mirror:simultaneous} extends strictly on these
new units to a functor
\begin{equation}
 J:\overline{\mathcal D}\longrightarrow\overline{\mathcal B}.
 \label{per:functor}
\end{equation}
Extend $\mathcal M_0$ by the same unit convention.  For strict arrows,
$J$ induces cohomology isomorphisms.  Each $w_X:X\to X^+$ maps to a
cohomological isomorphism, and its action on $\mathcal M_0$ is a
quasi-isomorphism by Proposition~\ref{cat:continuation}.

We use strictly unital right $A_\infty$ modules and their dg category
of premorphisms.  All degrees in the rest of this section are modulo
four unless explicitly stated otherwise.  The parity signs are the
usual ones, since the period is even.  Write $h_x$ for the right
representable at an object $x$.

\begin{lemma}[Periodic module calculus]
\label{per:module-calculus}
For an uncurved strictly unital $A_\infty$ category over a field,
the dg Hom of strictly unital modules computes
derived Hom after objectwise quasi-isomorphisms are inverted.  Direct
sums are coproducts, and
\[
 \Hom^*(h_x,N)=H^*(N(x)).
\]
The normalized bar tensor is exact on triangles and preserves
objectwise quasi-isomorphisms; tensoring with the diagonal gives the
identity.  These assertions hold with direct sums in bar length and
products in premorphism length in the period-four grading.
\end{lemma}

\begin{proof}
Filter a premorphism by its number of arrow inputs.  Its components
form a product, complete for the filtration by components of length
at least $n$.  The length-preserving part of the differential is the
internal Hom differential.  If the target module is objectwise
acyclic, that associated-graded Hom complex is acyclic: over a field
each acyclic value complex splits, and its contraction contracts Hom
from any complex.  This also applies to the normalized arrow inputs,
with the strict units divided out.  Starting at length zero, solve a
boundary equation recursively in increasing length.  At the next
length the already solved equations make the residual component an
internal cocycle, so it can be solved there.  Completeness of the
product gives the desired premorphism.  Thus Hom into an objectwise
acyclic target is acyclic.  In particular an objectwise acyclic
module is contractible, by applying this conclusion to its identity,
and quasi-isomorphisms have contractible cones.

Direct sums are formed on values and operations.  Their premorphism
Hom into a fixed target is the product of the individual Hom
complexes, giving the coproduct assertion.  The ordinary strictly
unital Yoneda homotopy evaluates a premorphism on the identity at
$x$ and contracts its remaining bar components by inserting that
identity.  Each component of this homotopy involves only a finite
word, so it works unchanged with periodic degrees.  It gives the
displayed Yoneda isomorphism, naturally for the cohomology actions.

The normalized bar tensor uses a direct sum of finite tensor words;
its differential does not increase word length.  If one factor is
objectwise acyclic, its length-preserving associated-graded
differential is acyclic.  A cycle involves only boundedly many
lengths.  Kill its highest-length component by an internal boundary
and continue downwards, proving that the bar tensor is acyclic.
Cones then prove invariance under quasi-isomorphisms and exactness
on triangles.  The diagonal contraction is the Yoneda bar
contraction: insert an identity in the end representable slot.
Strict unitality removes the other higher-unit terms.  All these
identities are finite at each word, so none requires convergence
with respect to integer cohomological degree.
\end{proof}

\subsection{Localization by later copies}

Define a right module over the full folded line model by coinduction:
\begin{equation}
 \widetilde{\mathcal M}(c)
   =\RHom_{\overline{\mathcal D}}
       \bigl(J^*h_c^{\overline{\mathcal B}},\mathcal M_0\bigr).
 \label{per:coinduction}
\end{equation}
The restricted Yoneda functor takes $c$ covariantly to a
$\overline{\mathcal D}$-module; dg Hom into $\mathcal M_0$ makes
\eqref{per:coinduction} a right module.  The unit of $Jx$ gives the
canonical Yoneda map
\[
 u_x:h_x^{\overline{\mathcal D}}\longrightarrow
                  J^*h_{Jx}^{\overline{\mathcal B}}.
\]
Precomposition defines an evaluation comparison
\begin{equation}
 H^*\bigl(\widetilde{\mathcal M}(Jx)\bigr)
       \longrightarrow H^*\bigl(\mathcal M_0(x)\bigr).
 \label{per:evaluation-map}
\end{equation}

\begin{proposition}[Evaluation after localization]
\label{per:localization}
The comparison \eqref{per:evaluation-map} is an isomorphism for every
graph label $x$.  It respects the cohomology actions of graph arrows
under $J$.  In particular $\widetilde{\mathcal M}$ has finite-dimensional
total cohomology at every object of $\overline{\mathcal B}$.
\end{proposition}

\begin{proof}
Fix $x_0=x$ and put $x_{j+1}=x_j^+$.  The layer ordering makes $x_j$
strictly later than any fixed label once $j$ is large.  Form the
sequential telescope of right representables, defined by the triangle
\begin{equation}
 \bigoplus_{j\ge0}h_{x_j}
  \xrightarrow{\,1-\mathrm{shift}\,}
 \bigoplus_{j\ge0}h_{x_j}
  \longrightarrow T_x\longrightarrow
 \left(\bigoplus_{j\ge0}h_{x_j}\right)[1],
 \qquad
 T_x=\hocolim_j h_{x_j}.
 \label{per:telescope}
\end{equation}
The shift maps are induced by $w_{x_j}$.  In
$H^0(\overline{\mathcal B})$, let
$b_j:Jx_j\to Jx$ be the inverse of the composite from $Jx$ to $Jx_j$,
with $b_0=\id$.  The Yoneda classes defined by $b_j$ give compatible
maps
$h_{x_j}\to J^*h_{Jx}^{\overline{\mathcal B}}$: their composites with
$w_{x_j}$ agree on cohomology.  Exactness of Hom applied to
\eqref{per:telescope} therefore supplies a map
\begin{equation}
 T_x\longrightarrow J^*h_{Jx}^{\overline{\mathcal B}}
 \label{per:telescope-comparison}
\end{equation}
whose restriction to $h_x$ is $u_x$.

Evaluate this map at a label $a$.  Cohomology of the telescope is the
direct limit of the cohomology of its value complexes.  For all large
$j$, the arrows $a\to x_j$ are strict, and hence $J$ induces an
isomorphism
\[
 H^*\overline{\mathcal D}(a,x_j)
    \xrightarrow{\ \sim\ }
 H^*\overline{\mathcal B}(Ja,Jx_j).
\]
Composing with $b_j$ identifies the right side with
$H^*\overline{\mathcal B}(Ja,Jx)$.  These identifications commute with
the transition maps.  Thus \eqref{per:telescope-comparison} is an
objectwise quasi-isomorphism.

Apply derived Hom into $\mathcal M_0$.  By
Lemma~\ref{per:module-calculus}, the two direct sums in
\eqref{per:telescope} become products of the complexes
$\mathcal M_0(x_j)$.  Products of vector spaces are exact, so their
cohomology is the product of the value cohomologies.  The difference
map on these products is surjective: all inverse-system transition
maps
\[
 H^*\mathcal M_0(x_{j+1})\longrightarrow H^*\mathcal M_0(x_j)
\]
are isomorphisms, and one solves the difference equation recursively.
Consequently there is no derived inverse-limit term.  The cohomology
of $\RHom(T_x,\mathcal M_0)$ is the inverse limit of these groups,
and projection to its zeroth factor is an isomorphism.  In view of
\eqref{per:telescope-comparison}, it is exactly
\eqref{per:evaluation-map}.

For naturality, an arrow $x\to y$ gives two composites from $h_x$ to
$J^*h_{Jy}$.  By Yoneda both represent its image under $J$, since the
maps $u_x,u_y$ represent the identity classes.  Precomposition
therefore makes \eqref{per:evaluation-map} natural for the
cohomology actions.  Finally the value complexes of $\mathcal M_0$
are finite-dimensional by Proposition~\ref{cat:module}.
\end{proof}

\subsection{A perfect object in the folded hull}

\begin{proposition}
\label{per:perfect}
The module $\widetilde{\mathcal M}$ is perfect over
$\overline{\mathcal B}$.  It is represented by an object
$\mathfrak E$ in the split-closed pretriangulated hull of the
direct-sum fold of an algebraic dg model of $\Perf(Y_\kappa)$.  For
each graph label $X$, its value cohomology identifies with the
periodic Hom from the corresponding ordinary line to $\mathfrak E$,
and these identifications respect cohomology compositions.
\end{proposition}

\begin{proof}
The integer-graded category $\mathcal B$ is smooth: its diagonal
bimodule is perfect, since its split-closed hull is
$\Perf(Y_\kappa)$ and $Y_\kappa$ is smooth projective.  This is the
smooth projective Morita comparison: the perfect kernel of the
geometric diagonal gives the diagonal bimodule; see
\cite{toen-vaquie}.  One may express
this by a finite presentation using cones, shifts and a retract of
exterior Yoneda cells.  Such a presentation remains a presentation of
the correct diagonal after folding.  Indeed for any two complexes
$K,L$,
\[
 \operatorname{Fold}(K\otimes L)^r
   =\bigoplus_{a+b\equiv r\;({\rm mod}\,4)}K^a\otimes L^b
   =\bigl(\operatorname{Fold}(K)\otimes
                      \operatorname{Fold}(L)\bigr)^r.
\]
The differentials and actions agree, including signs, because four
is even.  Therefore an exterior Yoneda cell folds to the exterior
cell of the folded representables, while the diagonal folds to the
actual diagonal of $\overline{\mathcal B}$.  Folding preserves
quasi-isomorphisms, finite cones and retracts.  Thus
$\overline{\mathcal B}$ is smooth in the periodic grading.

Tensor a split-perfect diagonal presentation with
$\widetilde{\mathcal M}$.  By Lemma~\ref{per:module-calculus}, this
recovers $\widetilde{\mathcal M}$, and every exterior cell becomes a
representable tensored with one of its value complexes.  By
Proposition~\ref{per:localization} that complex has finite-dimensional
cohomology.  Over a field it is therefore perfect as a periodic
complex: split off its acyclic part and take its four finite
cohomology spaces.  The tensor presentation proves perfectness.

For the geometric model, take bounded vector bundle complexes on
$Y_\kappa$, with dg Hom given by a fixed finite affine-cover \v Cech
complex and cup composition.  The ordinary comparison with the line
model is cohomologically fully faithful and split-generating.
Direct-sum folding preserves the first property because cohomology
commutes with direct sums.  It preserves the second because each
generation relation uses finitely many cones and a retract.  Taking
the split-closed pretriangulated hull gives the asserted object
$\mathfrak E$.  The representables $h_{G_k}$ and $h_C$ correspond to
the folds of $L_k'$ and $P_C$.  The evaluation and action assertions
are Proposition~\ref{per:localization} and Yoneda.
\end{proof}

The object $\mathfrak E$ is only asserted to belong to this periodic
hull.  We shall use its coherent periodic cohomology; no
integer-graded lift is needed.

\subsection{Identifying the polarization and section actions}

Let $\mathscr S_{\rm s}$ and
$\mathscr X_{\rm s}=\Proj\mathscr S_{\rm s}$ be as in
Proposition~\ref{theta:family}.

\begin{lemma}
\label{per:section-ring}
There are an isomorphism
$\phi:Y_\kappa\xrightarrow{\sim}\mathscr X_{\rm s}$ and a line
$B_0$ on $Y_\kappa$ such that, with
$A_1=\phi^*\cO_{\mathscr X_{\rm s}}(1)$,
\begin{equation}
 L_k'=B_0\otimes A_1^{-k},\qquad k\ge1.
 \label{per:line-progression}
\end{equation}
Under these identifications, the maps from
$\mathscr S_{{\rm s},i-j}$ to the cohomology arrows $G_i\to G_j$
are ordinary section multiplication.
\end{lemma}

\begin{proof}
For $j\ge1$, the degree-one arrow space identifies with
\[
 \mathscr S_{{\rm s},1}
    \xrightarrow{\sim}
 H^0\bigl(Y_\kappa,(L_{j+1}')^\vee\otimes L_j'\bigr).
\]
Here strict-arrow full faithfulness applies, and the quotient line
has ample product type $H_Y^{h_0}$, so its ordinary Ext occurs only
in degree zero.  The complete linear system is very ample: each
elliptic factor has degree $h_0\ge3$, and Picard-zero twists preserve
this property.  It defines a closed immersion
\[
 f_j:Y_\kappa\longrightarrow
        \mathbb P(\mathscr S_{{\rm s},1}^{\vee}).
\]

These morphisms are independent of $j$.  Indeed commutativity of
the section algebra, transported through compositions, says that
the evaluation functionals for two consecutive quotients satisfy
\[
 \lambda_j(a)\lambda_{j+1}(b)
       =\lambda_j(b)\lambda_{j+1}(a)
       \qquad(a,b\in\mathscr S_{{\rm s},1}).
\]
Both functionals are everywhere nonzero by global generation.
Their cross-product equations give the same projective quotient,
so $f_j=f_{j+1}$.  Write $f$ for this common closed immersion.

Every homogeneous relation in $\mathscr S_{\rm s}$ vanishes on
the image of $f$: express it as a sum of products of degree-one
sections, realize the products by successive line arrows, and
evaluate.  It follows that $f$ factors through
$\mathscr X_{\rm s}\subset
\mathbb P(\mathscr S_{{\rm s},1}^{\vee})$.  Both varieties have
dimension eight, and $\mathscr X_{\rm s}$ is smooth and integral
by Proposition~\ref{theta:family}.  Thus the closed image is all of
$\mathscr X_{\rm s}$; reducedness makes this an isomorphism, denoted
by $\phi$.

Each successive quotient line, with its evaluation, is the pullback
of the universal projective quotient $\cO(1)$.  This identifies all
the quotients with $A_1$, including their degree-one section maps.
Choose the line identifications successively and set
$B_0=L_1'\otimes A_1$.  This gives
\eqref{per:line-progression}.  Since $\mathscr S_{\rm s}$ is generated
in degree one, compatibility with degree-one multiplication gives
compatibility with every section arrow.
\end{proof}

\subsection{Coherent periodic cohomology}

Write $D_4(\operatorname{QCoh}Y_\kappa)$ for the derived category of
period-four complexes of quasi-coherent sheaves, with
quasi-isomorphisms of sheaf complexes inverted.  We allow its idempotent
completion.

\begin{lemma}[Finite \v Cech realization]
\label{per:coherent}
The folded perfect hull in Proposition~\ref{per:perfect} maps fully
faithfully on cohomology into the idempotent completion of
$D_4(\operatorname{QCoh}Y_\kappa)$.  Its objects have coherent
periodic cohomology sheaves.  For a vector bundle $P$ and an object
$E$ of this hull there is a natural, finite-column spectral sequence
\begin{equation}
 E_2^{a,b}
   =H^a\bigl(Y_\kappa,P^\vee\otimes\mathcal H^b(E)\bigr)
   \quad\Longrightarrow\quad \Hom^{a+b}(P,E),
 \qquad b\in\Z/4.
 \label{per:spectral-sequence}
\end{equation}
It is natural for ordinary bundle maps in $P$ and for idempotents
in $E$.
\end{lemma}

\begin{proof}
Fix a finite affine open cover of $Y_\kappa$ and use ordered
\v Cech cochains.  All its intersections are affine, since
$Y_\kappa$ is separated.  Send a bounded vector bundle complex $P$
to its sheaf \v Cech resolution, folded into period four.  A
\v Cech Hom cochain acts on these resolutions by the cup map, with
the usual total-complex signs.  Cup composition makes this a dg
construction.  It extends to finite twisted complexes and hence,
on cohomology, to their retracts in the idempotent-completed target.

We check full faithfulness first on folded bounded vector bundle
complexes $P,Q$.  Precompose a cup map by the augmentation
$P\to\check C(P)$.  The resulting maps from $P$ to
$\check C(Q)$ are precisely the global total \v Cech Hom
cochains of the original dg model.  They already compute derived
Hom.  To verify this last assertion, filter the sheaf \v Cech target
by its finitely many columns.  Each column is a finite sum of
terms $j_*K$, where $j:U\hookrightarrow Y_\kappa$ is an affine
intersection of the cover.  Such $j$ is an affine morphism:
its intersection with any affine open of the separated variety is
affine.  Therefore $j_*$ is exact on quasi-coherent sheaves;
restriction is exact as well.  Their ordinary adjunction consequently
passes to the periodic derived categories.  On $U$, the fold of the
bounded complex $P|_U$ consists of finite projectives and computes
outgoing Hom without further resolution.  Indeed its finite
ordinary-degree filtration before folding reduces this assertion
to the exactness of Hom from one projective module.  The comparison follows
for each column by adjunction and for their total by the finite
column triangles.  The sheaf \v Cech augmentation is a
quasi-isomorphism, also periodically: locally its augmented cover
complex contracts by inserting a cover open containing the testing
neighborhood.  These facts prove full faithfulness on $P,Q$.
An exact functor fully faithful on those objects is fully faithful
on their finite cone hull, by the triangle Hom sequences, and then
on its retracts.

An ordinary folded perfect complex has coherent periodic cohomology.
The cohomology exact sequence of a cone preserves coherence, so this
holds for the finite twisted-complex hull.  For a retract, take the
images of the resulting idempotents on cohomology sheaves; these
images are again coherent.

For a pre-retract object choose its actual periodic sheaf-complex
representative $K$ obtained from the preceding construction.
Resolve it by the same sheaf \v Cech complex.  Exactness of
restriction and of the affine pushforwards shows that taking
internal cohomology first gives the \v Cech cochains of
$P^\vee\otimes\mathcal H^b(K)$.  Taking their cohomology gives
the left side of \eqref{per:spectral-sequence}.  The affine
derived-Hom argument just given identifies the total cohomology
with the right side.  There are only finitely many cover columns,
so the filtration and its abutment are finite in every residue
degree, although the internal grading is periodic.

The construction is functorial for chain maps and for roofs in
the derived category: a quasi-isomorphism induces an isomorphism
already on the internal-cohomology page.  Thus an idempotent acts
on the spectral sequence and its finite abutment filtration.
Taking its images proves \eqref{per:spectral-sequence} for a
retract.  Cup maps from ordinary bundle morphisms make the asserted
source naturality explicit.
\end{proof}

Apply this lemma to $\mathfrak E$ and abbreviate
$\mathcal H^b=\mathcal H^b(\mathfrak E)$.

\begin{proposition}[Coherent section-module tails]
\label{per:finite-generation}
For every $r\in\Z/4$, the module $V_{\rm s}^r$ is finitely generated
over $\mathscr S_{\rm s}$ and has associated sheaf
\begin{equation}
 \widetilde{V_{\rm s}^r}
   \cong\phi_*\bigl(B_0^\vee\otimes\mathcal H^r\bigr).
 \label{per:associated-sheaf}
\end{equation}
For all sufficiently large $j$, simultaneously for the four values
of $r$, its degree-$j$ piece is
\begin{equation}
 (V_{\rm s}^r)_j
  \cong H^0\bigl(Y_\kappa,
       A_1^j\otimes B_0^\vee\otimes\mathcal H^r\bigr),
 \label{per:tails}
\end{equation}
compatibly with multiplication by $\mathscr S_{\rm s}$.
\end{proposition}

\begin{proof}
By Propositions~\ref{per:localization} and~\ref{per:perfect},
$(V_{\rm s}^r)_j$ is $\Hom^r(L_j',\mathfrak E)$, with the given
section actions.  Use \eqref{per:line-progression} in
\eqref{per:spectral-sequence}.  Serre vanishing kills all terms
with $a>0$ for sufficiently large $j$, uniformly over the four
coherent sheaves $B_0^\vee\otimes\mathcal H^b$.  This gives
\eqref{per:tails}.  Naturality in the line source, together with
Lemma~\ref{per:section-ring}, identifies the multiplication actions.

The ring $\mathscr S_{\rm s}$ is the section ring of
$\cO_{\mathscr X_{\rm s}}(1)$, as established in the theta
construction.  Serre's section-module theorem therefore makes
each displayed tail finitely generated and identifies its Proj
sheaf.  The finitely many earlier pieces of $V_{\rm s}^r$ are
finite-dimensional.  Adjoining bases for them to tail generators
proves finite generation of the whole module and gives
\eqref{per:associated-sheaf}; see \cite{serre-proj}.
\end{proof}

\subsection{Euler pairings detect a nonpolarization class}

Fix a prime $\ell$ and use rational $\ell$-adic cycle cohomology,
with the Tate twists implicit in Chern characters.  Put
\begin{equation}
 \begin{split}
 l&=c_1\bigl(\cO_{\mathscr X_{\rm s}}(1)\bigr),\\
 \zeta_{\rm s}
   &=\sum_{r\in\Z/4}(-1)^r
                  \ch\bigl(\widetilde{V_{\rm s}^r}\bigr),\\
 \eta_{\mathfrak E}
   &=\sum_{b\in\Z/4}(-1)^b\ch(\mathcal H^b).
 \end{split}
 \label{per:characters}
\end{equation}
These are rational algebraic Chern-character classes.  By
\eqref{per:associated-sheaf},
\begin{equation}
 \phi^*\zeta_{\rm s}
       =e^{-c_1(B_0)}\eta_{\mathfrak E}.
 \label{per:character-comparison}
\end{equation}

For a vector bundle $P$, define the periodic Euler pairing by
\[
 \chi(P,\mathfrak E)
    =\sum_{r\in\Z/4}(-1)^r
                       \dim_\kappa\Hom^r(P,\mathfrak E).
\]
The groups are finite by \eqref{per:spectral-sequence} and properness.
Since the period is even and the column filtration finite, the
alternating sum on that spectral sequence is unchanged from its
second page to its abutment.  Ordinary Riemann--Roch, and the trivial
Todd class of an abelian variety, give
\begin{equation}
 \chi(P,\mathfrak E)
    =\int_{Y_\kappa}\ch(P^\vee)\,\eta_{\mathfrak E}.
 \label{per:rr}
\end{equation}
Both sides are additive in triangles in the ordinary perfect source,
so the identity also holds for folds of ordinary perfect complexes,
using a bounded vector bundle resolution.

\begin{proposition}[Nonpolarization detection]
\label{per:nonpolarization}
The class $\zeta_{\rm s}$ satisfies
\begin{equation}
 \zeta_{\rm s}\notin
       \operatorname{span}_{\Q_\ell}\{1,l,l^2,\ldots,l^8\}.
 \label{per:nonpolarization-equation}
\end{equation}
\end{proposition}

\begin{proof}
Let $P_k=L_k'$ and retain $P_C$ for the ordinary line corresponding
to $C$.  Lemma~\ref{cat:euler} and
Proposition~\ref{per:localization} identify
\[
 \chi(P_k,\mathfrak E)=\nu\,g_k\cdot\alpha,
 \qquad
 \chi(P_C,\mathfrak E)=\nu\,c\cdot\alpha.
\]
Here $\alpha=\alpha_{\rm ex}+g_1+g_2+g_3$,
$g_k\cdot\alpha_{\rm ex}=0$, and
$c\cdot\alpha_{\rm ex}\ne0$.

Ordinary graph Euler pairings also give
\begin{equation}
 \chi(P_k,P_j)=g_k\cdot g_j,
 \qquad
 \chi(P_C,P_j)=c\cdot g_j.
 \label{per:graph-pairings}
\end{equation}
In the available strict-arrow direction these follow from full
faithfulness and the ordinary intersection-sign formula.  Reversing
the direction preserves Euler pairings by Serre duality in the even
dimension eight, just as it preserves the middle intersection form.
For equal labels both self-Euler numbers are zero.

Suppose $\zeta_{\rm s}$ belongs to the span in
\eqref{per:nonpolarization-equation}.  Equations
\eqref{per:line-progression}, \eqref{per:character-comparison},
\eqref{per:rr}, and \eqref{per:graph-pairings} imply, for every
$k\ge1$,
\begin{align}
 \int_{\mathscr X_{\rm s}}e^{kl}\zeta_{\rm s}
   &=\chi(P_k,\mathfrak E)
     =\nu\sum_{j=1}^3\chi(P_k,P_j)\notag\\
   &=\nu\int_{\mathscr X_{\rm s}}
                         e^{kl}\sum_{j=1}^3e^{-jl}.
 \label{per:scalar-test}
\end{align}
These evaluations determine every coefficient of a polynomial in
$l$.  Explicitly, if $z=\sum_{a=0}^8z_al^a$, then
\[
 \int_{\mathscr X_{\rm s}}e^{kl}z
    =\left(\int_{\mathscr X_{\rm s}}l^8\right)
       \sum_{a=0}^8\frac{z_a}{(8-a)!}k^{8-a}.
\]
The leading intersection factor is nonzero because $l$ is ample.
Equality for all positive integers $k$ in \eqref{per:scalar-test}
therefore forces
\[
 \zeta_{\rm s}=\nu\sum_{j=1}^3e^{-jl},
 \qquad
 \eta_{\mathfrak E}=\nu\sum_{j=1}^3\ch(P_j).
\]
Testing the second equality by $\ch(P_C^\vee)$ in
\eqref{per:rr} would give
\[
 \chi(P_C,\mathfrak E)=\nu\sum_{j=1}^3\chi(P_C,P_j).
\]
The two sides differ, by the Euler and intersection identities,
by $\nu\,c\cdot\alpha_{\rm ex}\ne0$.  This contradiction proves
\eqref{per:nonpolarization-equation}.
\end{proof}
